% SEGMENT OLP-0351-S01
% Part: lambda-calculus
% Chapter: introduction
% Section: basic-pr-lambda

\documentclass[../../../include/open-logic-section]{subfiles}

\begin{document}

\olfileid{lam}{int}{bas}
\olsection{அடிப்படை முதன்மை மீள்வரையறைச் சார்புகள் \usetoken{S}{lambda definable}}

\begin{lem}
சார்புகள் $\Zero$, $\Succ$, $\Proj{n}{i}$ ஆகியவை !!{lambda definable}.
\end{lem}

\begin{proof}
$\Zero$ என்பது $\lambd[x][\lambd[y][y]]$ மட்டுமே.

அடுத்தெண் சார்பு~$\Succ$ ஆனது
$\fn{Succ}(u) = \lambd[x][\lambd[y][x(uxy)]]$ ஆல் வரையறுக்கப்படுகிறது.
இது ஏன் செயல்படுகிறது என்று நீங்கள் சிந்திக்க வேண்டும். மீள்செயல்படுத்தியாகக்
கருதப்படும் ஒவ்வோர் எண்ணுரு $\num{n}$ மற்றும் ஒவ்வொரு சார்பு~$f$ க்கும்,
$\fn{Succ}(\num{n},f)$ ஆனது உள்ளீடு~$y$ மீது,~$y$ இலிருந்து தொடங்கி $f$ ஐ
$n$ முறை செயல்படுத்திவிட்டு, பின்னர் அதை மேலும் ஒருமுறை செயல்படுத்தும் சார்பு.

வீழ்த்தல்களைப் பற்றி வேறெதுவும் கூறத் தேவையில்லை:
$\fn{Proj}^n_i(x_0, \dots, x_{n-1}) = x_i$. வேறு சொற்களில், நமது
மரபுகளின்படி $\fn{Proj}^n_i$ என்பது
$\lambd[x_0][\dots \lambd[x_{n-1}][x_i]]$ என்ற லாம்டா சொல்.
\end{proof}

\end{document}
